\chapter*{Course Map}
\phantomsection%
\addcontentsline{toc}{chapter}{Course Map}
\markboth{Course Map}{}

Start with the command-line environment, then follow the operating system from
\textbf{processes} to \textbf{CPU time}, \textbf{shared resources}, \textbf{memory},
and finally \textbf{persistent storage}. Each numbered chapter is followed by the
lab that puts that chapter's ideas into practice, and the capstone brings the five
stages together.

\begin{utpfigure}
\newcommand{\coursemapcard}[3]{%
  \parbox[c][1.52cm][t]{3.35cm}{\raggedright
    {\scriptsize\bfseries\color{UTPTeal}\strut #1}\par
    \vspace{0.02em}%
    {\bfseries\strut #2}\par
    \vspace{0.08em}%
    {\scriptsize\color{UTPMuted!80!UTPNavy}\strut #3\strut}%
  }%
}
\begin{tikzpicture}[
  node distance=0.58cm and 0.66cm,
  card/.style={draw=UTPBlue,fill=UTPPaleBlue,rounded corners=2mm,
    minimum width=3.99cm,minimum height=1.88cm,align=left,
    inner xsep=3.2mm,inner ysep=1.6mm,font=\small\color{UTPNavy}},
  finalcard/.style={card,draw=UTPTeal,fill=UTPPaleTeal},
  arr/.style={-{Stealth[length=2.25mm,width=1.55mm]},line width=0.9pt,
    draw=UTPTeal,shorten >=0.8pt,shorten <=0.8pt}
]
\node[card] (a) {\coursemapcard{STAGE 01}{Orientation}{Linux, the terminal, and the OS}};
\node[card,right=of a] (b) {\coursemapcard{STAGE 02}{Processes}{Processes, files, and the shell}};
\node[card,right=of b] (c) {\coursemapcard{STAGE 03}{CPU and concurrency}{Scheduling, IPC, and deadlock}};
\node[card,below=of c] (d) {\coursemapcard{STAGE 04}{Memory}{Allocation, paging, and replacement}};
\node[card,left=of d] (e) {\coursemapcard{STAGE 05}{Storage}{Directories, records, and blocks}};
\node[finalcard,left=of e] (f) {\coursemapcard{CAPSTONE}{Exam map}{Trace state across the course}};
\draw[arr] (a) -- (b);
\draw[arr] (b) -- (c);
\draw[arr] (c) -- (d);
\draw[arr] (d) -- (e);
\draw[arr] (e) -- (f);
\end{tikzpicture}
\end{utpfigure}

\begin{utptable}
\begin{tabularx}{\linewidth}{@{}>{\bfseries\leavevmode\color{UTPNavy}\centering\arraybackslash}m{0.08\linewidth}m{0.31\linewidth}X@{}}
\toprule
\rowcolor{UTPPaleBlue!92}
\utpth{Stage} & \utpth{Focus} & \utpth{Key question} \\
\midrule
1 & Linux, the toolchain, and the OS's role & Where does the OS sit between a program and the hardware? \\
2 & Processes, files, and shell automation & How is live OS state inspected and reported reliably? \\
3 & Scheduling, IPC, threads, and deadlock & Why does a policy or a lock change the outcome? \\
4 & Contiguous allocation, paging, and replacement & Which placement or victim choice does each policy make, and why? \\
5 & Directories, records, and block allocation & How do the names an application sees map to storage below? \\
\bottomrule
\end{tabularx}
\end{utptable}

\begin{keyidea}
Learn the idea first, then test it. In every lab, keep three questions visible:
\textbf{What state exists? What event changes it? What rule decides the next state?}
\end{keyidea}

\section*{Chapters, labs, and checkpoints}

\begin{utptable}
\begin{tabularx}{\linewidth}{@{}>{\bfseries\leavevmode\color{UTPNavy}\centering\arraybackslash}m{0.08\linewidth}>{\centering\arraybackslash}m{0.20\linewidth}X@{}}
\toprule
\rowcolor{UTPPaleBlue!92}
\utpth{Stage} & \utpth{Where} & \utpth{Checkpoint} \\
\midrule
1 & Setup and the OS primer & Open, edit, compile, run, and explain where the OS sits between a program and hardware. \\
2 & Chapters 1--3, Labs 1--3 & Inspect live OS state and build small tools that report it reliably. \\
3 & Chapters 4--6, Labs 4--6 & Trace state changes and explain why a policy or lock changes the outcome. \\
4 & Chapters 7--8, Labs 7--8 & Compare policies on the same workload and justify every placement or victim choice. \\
5 & Chapters 9--10, Labs 9--10 & Separate the names and records seen by applications from the storage structures managed below them. \\
\bottomrule
\end{tabularx}
\end{utptable}
